% Generated draft: requires mathematical and source review.
\documentclass[11pt]{article}
\usepackage[margin=25mm]{geometry}
\usepackage{fontspec}
\setmainfont{Latin Modern Roman}
\usepackage{amsmath,amssymb,mathtools,unicode-math}
\setmathfont{Latin Modern Math}
\usepackage{enumitem,needspace,xurl,hyperref}
\setlength{\parindent}{0pt}
\setlength{\parskip}{5pt}
\setlength{\emergencystretch}{3em}
\newcommand{\N}{\mathbb{N}}
\newcommand{\Z}{\mathbb{Z}}
\newcommand{\Q}{\mathbb{Q}}
\newcommand{\R}{\mathbb{R}}
\newcommand{\C}{\mathbb{C}}
\newcommand{\F}{\mathbb{F}}
\newcommand{\A}{\mathbb{A}}
\newcommand{\PP}{\mathbb{P}}
\newcommand{\E}{\mathbb{E}}
\newcommand{\eps}{\varepsilon}
\newcommand{\dd}{\,\mathrm{d}}
\begin{document}
\section*{UnsolvedMath ID: 701743}
\textbf{Source problem code:} catalogue entry
\subsection{Definitions and mathematical statement}
Let \(F\) be a free group; the source imposes no finite-rank hypothesis. In the usual free-group terminology, an element \(w\in F\) is \emph{primitive} if it belongs to a free basis of \(F\). The supplied record instead says that it belongs to a ``minimal generating system''; the meaning of ``minimal'' is not specified, so the free-basis interpretation requires confirmation.

Following the supplied wording, \(w\) is \emph{almost primitive} if, for every finitely generated subgroup \(H<F\) with \(w\in H\), the element \(w\) is primitive in \(H\). The supplied definition does not additionally require that \(w\) be nonprimitive in \(F\).

For a subgroup \(H\le F\), put
\[
 S_H(w)=\{n\in\mathbb Z_{\ge1}:w^n\in H\}.
\]
The element \(w\) is \emph{tame almost primitive} if, for every subgroup \(H\le F\) for which \(S_H(w)\ne\varnothing\), setting \(\alpha=\min S_H(w)\) gives
\[
 w^\alpha\text{ is primitive in }H
 \quad\text{or}\quad
 [F:H]=\alpha.
\]
Here \([F:H]\) is the subgroup index. This definition quantifies over all subgroups, without a finite-generation or properness restriction. If \(S_H(w)=\varnothing\), it imposes no condition for that subgroup. No separate nonprimitivity condition is supplied.

For integers \(g\ge1\) and \(p\ge2\), let
\[
\begin{aligned}
 F_{2g}&=\langle a_1,b_1,\ldots,a_g,b_g\rangle_{\mathrm{free}},
 &u_g&=\prod_{i=1}^{g}[a_i,b_i],\\
 F_p&=\langle c_1,\ldots,c_p\rangle_{\mathrm{free}},
 &v_p&=\prod_{j=1}^{p}c_j^2,\\
 F_{2g+p}&=\langle a_1,b_1,\ldots,a_g,b_g,c_1,\ldots,c_p\rangle_{\mathrm{free}}.
\end{aligned}
\]
Products are taken in increasing index order. Brackets denote group commutators; the supplied record does not specify the commutator convention. The source attributes to Fine, Moldenhauer, Rosenberger, and Wienke (2021) the results that \(u_g\) is tame almost primitive in \(F_{2g}\) and \(v_p\) is tame almost primitive in \(F_p\).

\textbf{Question (G. Rosenberger and L. Wienke, Kourovka Notebook 20.85).} Are there tame almost primitive elements in free groups other than the displayed elements \(u_g\), \(v_p\), and their products \(u_gv_p\) in \(F_{2g+p}\)?

\emph{Unresolved scope.} The source does not specify the equivalence relation intended by ``other than'': literal inequality, changes of free basis, isomorphisms of ambient free groups, conjugacy, or inversion may give different formulations. It also does not specify whether primitive elements should be excluded. Thus a precise inequivalence condition for an additional example cannot be supplied from this record alone. The product family is named in the question, but the preceding attribution explicitly asserts tameness only for the separate families \(u_g\) and \(v_p\).
\subsection{Short English statement}
The question asks whether free groups contain further tame almost primitive elements beyond the listed commutator, square, and combined product families. Tameness concerns the least positive power entering each subgroup: that power must be primitive there unless its exponent equals the subgroup index. The intended meaning of a further example, and whether primitive elements are excluded, remain unspecified.
\subsection{Sources}
\begin{itemize}
\item Local UnsolvedMath registry: permanent ID \(701743\), exact source problem code \texttt{catalogue entry}, title \emph{Kourovka Notebook Problem 20.85}. The supplied registry labels the problem open; that label is supplied source data, not an independently verified mathematical or literature assessment.
\item \emph{The Kourovka Notebook (Unsolved Problems in Group Theory)}, Problem 20.85, proposed by G. Rosenberger and L. Wienke in 2022. The supplied background identifies the 20th issue (2022) and the 21st issue edition (2026), \texttt{arXiv:1401.0300v48}. The supplied collection description names E. I. Khukhro and V. D. Mazurov as editors. These bibliographic details have not been independently checked.
\item B. Fine, A. Moldenhauer, G. Rosenberger, and L. Wienke, \emph{Topics in Infinite Group Theory: Nielsen Methods, Covering Spaces, and Hyperbolic Groups}, De Gruyter, Berlin, 2021. The supplied statement attributes the tameness results for \(u_g\) and \(v_p\) to this book; no page or theorem number is supplied.
\end{itemize}
\end{document}
